\scanpage{18}

\mspage{8}

a $SP_M=P_M$ i.e. $SF_M=F_M$ i.e. $\Lambda^*\simeq\{\pm1\}$, \uline{et} $\Lambda=\Lambda_0$ i.e. $\Lambda\simeq\mathbb{Z}$ ou $\Lambda\simeq\mathbb{Z}/n\mathbb{Z}$, $n\geq3$.

Notons que la structure connue des $(\mathbb{Z}/n\mathbb{Z})^*$ montre alors que les seuls cas sont ceux \deleted{où \unsure{on a}} \inserted{de}

\[
\Lambda=\mathbb{Z},\quad\Lambda=\mathbb{Z}/n\mathbb{Z}\quad\text{avec }n\in\{3,4,6\}.\tag*{(42)}
\]

Revenons au cas général, \deleted{et regardant} et notons que $S\widetilde{F}_M$, donc aussi $SP_M$, ne \deleted{dépend} change pas si on remplace la jauge $e_M$ par $-e_M$. Ainsi le groupe

\[
(\mathrm{GL}_{\Lambda}M)^{\pm}\overset{\mathrm{def}}{=}\{g\in\mathrm{GL}_{\Lambda}(\deleted{\unsure{N}}M)\mid\det g\in\{\pm1\}\}\tag*{(43)}
\]

opère sur le $2$-polyèdre (généralisé) $SF_M$. Nous allons introduire sur celui-ci une \uline{orientation}, qui elle dépend du choix de $e_M$, et est remplacée par l'orientation opposée, quand on remplace $e_M$ par $-e_M$. Pour ceci, revenons à l'application $\alpha_f$ (31)

\[
\alpha_f:\underline{\omega}_f\longrightarrow(\det M)^*
\]

% archive sheet number: 17
